% Section 4 -- Evaluation protocol.
% Contribution C6 lives here, together with the three pre-registered designs that could not answer their own questions.
\section{Evaluation Protocol}
\label{sec:protocol}

The results in this paper were produced under a protocol fixed before any of them existed. We describe it in detail
because two of its features---the amendment mechanism and the over-performance stop rule---did more work than any
single experiment, and because three of our own pre-registered designs turned out to be unable to answer the questions
they were written for. Reporting that is the point of this section.

\subsection{Pre-registration}
\label{sec:prereg}

Before any GPU work, a thresholds document was committed to the repository specifying, for each task, what would
count as a pass, what would count as a failure, what would be reported as marginal, and---separately, so that outcomes
could not be reinterpreted afterwards---what each verdict would license the paper to claim. The document was written
against the state of knowledge at that time, and it was wrong in instructive ways. It is preserved in the repository
with every superseded sentence struck through rather than deleted, so a reader can audit what was pre-registered
against what was concluded.

\subsection{Controls applied to every result}
\label{sec:controls}

Four controls are applied by default rather than on suspicion.

\textbf{Length stratification.} Length as a confound in language-model evaluation is documented: response length
biases uncertainty-quantification benchmarks \cite{lengthconfound2025a}, and at least one probe study length-matches
its conditions for exactly this reason \cite{evalawareness2026}. Prompt length in activation-probe safety evaluation is nonetheless
not routinely controlled, and our contribution here is to quantify what that costs. It is large on two independent
benchmarks used here. On AILuminate, character length alone separates harmful prompts from
the benign controls at AUC $0.95$, and probe cells that appeared to reach $0.85$--$0.92$ on encoded prompts fall to
$0.49$--$0.65$ once it is controlled (Section~\ref{sec:obfuscation}). On the harmful-content evaluation set, length
alone scores $0.903$ unstratified and $0.578$ length-stratified (Table~\ref{tab:af_lexical}) --- it is the single
strongest cheap feature available, and stratification is what removes it.
Every AUC we report is therefore accompanied by a length-stratified counterpart, written \aucstar{}: character length
is split into five quantile bins within the compared groups, an AUC is computed inside each bin, and the results are
pooled weighted by $n_{\text{pos}} \cdot n_{\text{neg}}$, with bins holding fewer than five of either class dropped
and reported. Where a threshold in this paper is stated as length-controlled, \aucstar{} is the number that must clear
it.

\textbf{Confidence intervals.} Every headline effect is reported with a $95\,\%$ interval in
Appendix~\ref{app:ci}: a nonparametric bootstrap over cases---$2{,}000$ resamples at a recorded seed, percentile
method---computed from the per-case scores committed beside every aggregate. A difference between two reads of the
same prompts is resampled paired rather than treating the arms as independent samples. The per-case scores are
first-class results in the repository, and a test fails any results directory that reports an AUC without them,
because an aggregate whose scores are gone can be quoted but never re-interrogated.

\textbf{Topic-matched negatives.} Every harmful-content result is reported against an on-topic benign set as well as
an off-topic one (Section~\ref{sec:data}), because the gap between the two is the measurement of interest rather than
a robustness check.

\textbf{Refusal separation.} A probe that fires on exactly the prompts the model was going to refuse anyway adds
nothing to a deployed system. For every task we generate the model's own continuation on a seeded subsample, classify
it as a refusal or not, and report probe detection conditioned on that. The marginal value of a probe over the
model's existing behaviour is then visible rather than assumed.

\textbf{Lexical baselines.} Every result is reported beside the best cheap text classifier we can construct for the
same contrast: for harmful-content detection, tuned keyword and request-shape features (Table~\ref{tab:af_lexical});
for policy compliance, a second-person-pronoun count and, later, a bag-of-words model. This is not a formality. In one
of the three tasks the lexical baseline matched the probe, and the finding that follows from it
(Section~\ref{sec:learn}) is the paper's main negative result.

\subsection{Leakage and grouping}
\label{sec:leakage}

Splits are checked for shared content, not merely for shared examples. The agent-injection data is built by crossing a
small pool of injection strings with a small pool of user instructions, so a split that is disjoint at the case level
can still place every test-half injection string in the training half; we report both that split and a grouped split
in which whole injection strings are assigned to one side (Section~\ref{sec:read}). A related point applies to
cross-validation on data containing near-duplicates: leave-one-out is invalid there, because a held-out item's twin
sits in the training fold carrying the opposite label and the model reliably anti-predicts. We measured this at AUC
$0.019$---near-perfect separation with the sign inverted---on a set built from matched pairs. The correct procedure,
which we adopt, is leave-pair-out cross-validation scored against a within-pair-swap null.

\subsection{Amendments}
\label{sec:amendments}

A pre-registered interpretation can fail in two distinct ways, and we record them distinctly.

\emph{Type A, mechanism reassigned}: the threshold is met and the effect is real, but the pre-registered cause is
refuted by a control. The claim survives with a different mechanism, and the original wording is struck through beside
the replacement.

\emph{Type B, evidence disqualified}: the threshold is met, but the evaluation data cannot bear the claim. The
constructive sentence is withheld, and no amount of additional analysis on the same data can restore it.

Both occurred. Section~\ref{sec:template} is a type-A amendment: harmful-content probing passed its criteria, but the
control showed the improvement came from somewhere other than where the design predicted. The pre-registered mechanism
was the paired data construction---on-topic benign twins cancelling topic in the difference of class means. The
ablation that holds the estimator and the positives fixed and swaps the twins for the off-topic negatives moved AUC
against on-topic benign prompts by $-0.003$ to $+0.018$, while reading the identical probes through the chat template
moved length-stratified AUC by $+0.30$ to $+0.54$; the mechanism was reassigned to extraction convention, and the
original sentence is struck through beside the replacement in the thresholds document. Section~\ref{sec:learn} is a
type-B amendment: policy-compliance probing passed its criteria on data that a regular expression can also solve.

\subsection{The over-performance stop rule}
\label{sec:stoprule}

The rule that caught the most is the one that treats a good result as a warning. A result that \emph{beats} its
pre-registered threshold stops the run and is not written up as a pass until contamination, leakage and cheap-baseline
checks have been run, when any of the following hold: it exceeds the threshold by at least $0.05$ AUC; any AUC is at
least $0.95$; or it exceeds the claims of prior work on the same benchmark.

The rule exists because, in the history of this codebase, an impressive number has repeatedly measured something other
than the claim: evaluation data that had silently fallen back to a synthetic generator, a training set that contained
the very encodings being tested, and prompt length. All three arms of the present study triggered the rule. One
overshoot survived the checks unchanged, one survived but changed mechanism, and one did not survive at all. Without
the rule, all three would have been reported as successes.

We also apply a degenerate-score stop: constant scores, scores collapsing into a few repeated values, or any loader
exception halts the run. An exact $0\,\%$ or $100\,\%$ rate is not by itself degenerate---it is non-blocking when the
underlying score distribution is well spread---but every such instance is logged. Nineteen occurred here, almost all
from applying a shipped probe outside the distribution it was fitted on.

\subsection{Three designs that could not answer their own questions}
\label{sec:selfdefeating}

Three of our pre-registered designs were found, during execution, to be incapable of testing what they were written to
test. We report them because the failure modes generalise.

\emph{A comparison between two estimators that are the same estimator.} The design compared a paired construction
against an unpaired one on identical text, so that ``the method is the only thing that varies''. For complete,
aligned, equal-sized pairs the two are algebraically identical (Section~\ref{sec:learn}), and the two conditions agreed
bit for bit on every model. The comparison was replaced by an ablation that varies the training \emph{negatives}
instead, which is what the underlying hypothesis had actually been about.

\emph{A split that measured recognition rather than generalisation.} Described in Section~\ref{sec:leakage}; replaced
by a grouped split.

\emph{An evaluation set that a regular expression can solve.} Described in Section~\ref{sec:learn}; not replaceable,
and the reason a claim is withheld.

The general lesson is that a pre-registered design is a hypothesis about the measurement instrument as much as about
the phenomenon, and can be falsified in the same way. Recording that falsification is more useful than quietly
substituting a design that works.

\subsection{Provenance}
\label{sec:provenance}

Every results file embeds the git commit, the sha256 of the configuration and of every dataset file, the resolved
package versions, and the GPU. The runner validates that block before committing a result, and refuses to run at all
if the sha256 of any authored evaluation data differs from the version signed off by the author. Every table in this
manuscript is generated from those files by a script that aborts the build if a source is missing, so no number here
is typed by hand. The tables were generated from results commit \texttt{\resultscommit}.
